\subsection{Семинар «Делегирование в Kotlin» (2 часа)}

\subsubsection*{Теоретическое введение}

В языке Kotlin делегирование (delegation) реализуется с помощью ключевого слова \texttt{by}. Этот механизм позволяет классу делегировать реализацию интерфейса другому объекту, что способствует повторному использованию кода без наследования. Делегирование особенно полезно при реализации паттерна \textit{Decorator}, когда необходимо динамически расширять функциональность объекта.

Рассмотрим пример, в котором класс \texttt{LoggerPrinter} реализует интерфейс \texttt{Printer}, делегируя один метод и переопределяя другой:

\begin{lstlisting}[language=Kotlin, basicstyle=\ttfamily\small]
interface Printer {
    fun print(message: String)
    fun log(message: String)
}

class BasicPrinter : Printer {
    override fun print(message: String) {
        println("Printing: $message")
    }

    override fun log(message: String) {
        println("Log: $message")
    }
}

class LoggerPrinter(printer: Printer) : Printer by printer {
    // Метод log делегируется базовому printer
    // Метод print переопределяется
    override fun print(message: String) {
        println("[SECURE] Printing: $message")
    }
}
\end{lstlisting}

В этом примере:
\begin{itemize}
    \item метод \texttt{log} автоматически делегируется объекту \texttt{printer} (благодаря \texttt{by printer});
    \item метод \texttt{print} переопределяется для добавления дополнительной логики.
\end{itemize}

Такой подход позволяет гибко комбинировать поведения без дублирования кода и без необходимости строить сложные иерархии наследования, особенно в условиях отсутствия множественного наследования.

\medskip

Реализуйте несколько классов, демонстрирующих применение паттерна \textit{Decorator} с использованием механизма делегирования, встроенного в Kotlin. Обратите внимание: приведённые задачи трудно или невозможно решить с помощью наследования без дублирования кода или изобретения избыточно сложных иерархий (особенно в языках без множественного наследования).

\begin{enumerate}
    \item \textbf{Преподаватели.}  
    Осуществите ввод информации о преподавателях и выведите тех, у кого средний балл учащихся выше 4.  
    Преподаватель характеризуется ФИО, полом и средним баллом учащихся.  
    Дополнительно:
    \begin{itemize}
        \item если имеет категорию — дата последнего подтверждения и номер приказа;
        \item если имеет кандидатскую степень — тема диссертации, научное направление, дата защиты.
    \end{itemize}

    \item \textbf{Студенты.}  
    Введите информацию о студентах и выведите тех, у кого средний балл выше 4.  
    Студент характеризуется ФИО, полом и средним баллом.  
    Дополнительно:
    \begin{itemize}
        \item если платник — номер договора и признак отсутствия задолженности;
        \item если военнообязанный — номер приписного свидетельства/военного билета, название военкомата, признак прохождения службы.
    \end{itemize}

    \item \textbf{Остановочные пункты ЖД.}  
    Введите данные об остановочных пунктах пригородного железнодорожного транспорта и выведите те, где пассажиропоток превышает 1000 чел./сутки.  
    Пункт характеризуется названием, пассажиропотоком и расстоянием от вокзала.  
    Дополнительно:
    \begin{itemize}
        \item при наличии турникетного комплекса — количество павильонов и стоимость билета «на выход»;
        \item если рядом со станцией метро — название станции и расстояние до неё.
    \end{itemize}

    \item \textbf{ВУЗы.}  
    Введите информацию о вузах и выведите те, где обучается более 2000 студентов.  
    ВУЗ характеризуется названием, числом студентов, адресом и номером лицензии.  
    Дополнительно:
    \begin{itemize}
        \item при наличии аккредитации — номер и дата окончания;
        \item если выдаёт документы международного образца — стоимость выдачи и регион признания (например, «ЕС», «США», «весь мир»).
    \end{itemize}

    \item \textbf{Планшеты.}  
    Введите данные о планшетах и выведите те, что стоят менее 15\,000 рублей.  
    Планшет характеризуется названием, производителем и стоимостью.  
    Дополнительно:
    \begin{itemize}
        \item при поддержке SIM-карты — поддержка 3G и LTE;
        \item при поддержке Bluetooth — версия стандарта и возможность подключения клавиатуры.
    \end{itemize}

    \item \textbf{Смартфоны.}  
    Введите информацию о смартфонах и выведите те, что стоят менее 10\,000 рублей.  
    Смартфон характеризуется моделью, производителем и ценой.  
    Дополнительно:
    \begin{itemize}
        \item при поддержке Wi-Fi — стандарт и максимальная скорость;
        \item при поддержке геопозиционирования — поддержка GPS/ГЛОНАСС и количество каналов.
    \end{itemize}

    \item \textbf{Телевизоры.}  
    Введите данные о телевизорах и выведите те, что стоят менее 10\,000 рублей.  
    Телевизор характеризуется маркой, производителем и диагональю.  
    Дополнительно:
    \begin{itemize}
        \item при поддержке Smart TV — название платформы и предустановленные приложения;
        \item при поддержке многоканального звука — наличие оптического и поканальных аудиовыходов.
    \end{itemize}

    \item \textbf{Накопители.}  
    Введите информацию о накопителях и выведите те, у которых ёмкость менее 100 ГБ.  
    Накопитель характеризуется моделью, производителем и ёмкостью.  
    Дополнительно:
    \begin{itemize}
        \item при интерфейсе USB — версия стандарта и скорость передачи;
        \item если это жёсткий диск — количество поверхностей и физический размер (2.5", 3.5").
    \end{itemize}

    \item \textbf{Компьютеры.}  
    Введите данные о компьютерах и выведите те, что стоят менее 20\,000 рублей.  
    Компьютер характеризуется производителем, моделью и стоимостью.  
    Дополнительно:
    \begin{itemize}
        \item если моноблок — диагональ дисплея;
        \item при поддержке Wi-Fi — стандарт и максимальная скорость.
    \end{itemize}

    \item \textbf{Телеканалы.}  
    Введите информацию о телеканалах и выведите те, чья аудитория превышает 1\,000\,000 зрителей.  
    Телеканал характеризуется названием и размером аудитории.  
    Дополнительно:
    \begin{itemize}
        \item если эфирный — диапазон (МВ/ДМВ) и частота;
        \item если государственный — объём финансирования.
    \end{itemize}

    \item \textbf{Реки.}  
    Введите данные о реках и выведите те, длина которых превышает 20 км.  
    Река характеризуется названием и длиной.  
    Дополнительно:
    \begin{itemize}
        \item если судоходная — ширина и глубина русла;
        \item если используется для водоснабжения — суточный водозабор и название водопроводной станции.
    \end{itemize}

    \item \textbf{Памятники.}  
    Введите информацию о памятниках и выведите те, у которых рейтинг выше 2.5.  
    Памятник характеризуется названием, адресом и рейтингом.  
    Дополнительно:
    \begin{itemize}
        \item если посвящён одному человеку — ФИО и годы жизни;
        \item если связан с военными действиями — название и годы войны.
    \end{itemize}

    \item \textbf{ПИФы.}  
    Введите данные о паевых инвестиционных фондах и выведите те, что основаны до 2008 года.  
    ПИФ характеризуется названием управляющей компании, названием фонда и годом основания.  
    Дополнительно:
    \begin{itemize}
        \item для интервального ПИФа — даты начала и окончания следующего интервала;
        \item для ПИФа акций — эшелон акций (первый, второй, другие).
    \end{itemize}

    \item \textbf{Вклады.}  
    Введите информацию о банковских вкладах и выведите те, у которых доходность выше 9\%.  
    Вклад характеризуется названием, банком и процентной ставкой.  
    Дополнительно:
    \begin{itemize}
        \item если разрешены пополнения — за сколько дней до окончания можно вносить средства и минимальный взнос;
        \item если разрешено снятие — через сколько дней после открытия и размер неснижаемого остатка.
    \end{itemize}

    \item \textbf{Маршруты электричек.}  
    Введите данные о маршрутах и выведите те, длина которых превышает 50 км.  
    Маршрут характеризуется начальным и конечным пунктом, а также длиной.  
    Дополнительно:
    \begin{itemize}
        \item если экспресс — стоимость проезда и время в пути;
        \item если межобластной — названия соответствующих областей.
    \end{itemize}

    \item \textbf{Холодильники.}  
    Введите информацию о холодильниках и выведите те, высота которых превышает 1.5 м.  
    Холодильник характеризуется габаритами (высота, ширина, глубина), моделью и производителем.  
    Дополнительно:
    \begin{itemize}
        \item при наличии морозильной камеры — её габариты;
        \item при наличии зоны свежести — её габариты.
    \end{itemize}

    \item \textbf{Наушники.}  
    Введите данные о наушниках и выведите те, что стоят более 500 рублей.  
    Наушники характеризуются производителем, моделью, наличием активного шумоподавления и ценой.  
    Дополнительно:
    \begin{itemize}
        \item если беспроводные — радиус действия и признак цифрового сигнала;
        \item если вставные — количество сменных амбушюров и наличие регулятора громкости.
    \end{itemize}

    \item \textbf{Электронные книги.}  
    Введите информацию об электронных книгах и выведите те, что стоят более 3000 рублей.  
    Книга характеризуется моделью, производителем и стоимостью.  
    Дополнительно:
    \begin{itemize}
        \item при дисплее на электронных чернилах — цветность (ч/б или цветной) и наличие подсветки;
        \item при сенсорном интерфейсе — тип сенсора (ёмкостный, резистивный и т.д.).
    \end{itemize}

    \item \textbf{Устройства GPS.}  
    Введите данные об устройствах GPS и выведите те, что стоят более 5000 рублей.  
    Устройство характеризуется моделью, производителем, стоимостью и поддержкой ГЛОНАСС.  
    Дополнительно:
    \begin{itemize}
        \item если туристическое — водонепроницаемость, наличие магнитного компаса, тип батареи;
        \item если цветной дисплей — количество поддерживаемых цветов.
    \end{itemize}

    \item \textbf{Туристические палатки.}  
    Введите информацию о палатках и выведите те, у которых водостойкость дна превышает 7000 мм.  
    Палатка характеризуется производителем, моделью, водостойкостью тента и дна.  
    Дополнительно:
    \begin{itemize}
        \item при наличии внутренней палатки — её ширина и высота;
        \item если многоместная — количество мест и входов.
    \end{itemize}

    \item \textbf{Обогреватели.}  
    Введите данные об обогревателях и выведите те, что стоят более 1500 рублей.  
    Обогреватель характеризуется производителем, моделью, стоимостью и мощностью.  
    Дополнительно:
    \begin{itemize}
        \item если инфракрасный — тип (галогенный, карбоновый, кварцевый) и наличие автоматического вращения;
        \item если поддерживает программирование — возможность отсрочки старта и регулировки температуры.
    \end{itemize}

    \item \textbf{Автомобили.}  
    Введите информацию об автомобилях и выведите те, у которых пробег менее 50\,000 км.  
    Автомобиль характеризуется маркой, моделью, годом выпуска и пробегом.  
    Дополнительно:
    \begin{itemize}
        \item если электромобиль — ёмкость батареи и запас хода;
        \item если имеет систему помощи водителю — поддержка адаптивного круиз-контроля и автоматического торможения.
    \end{itemize}

    \item \textbf{Кофемашины.}  
    Введите данные о кофемашинах и выведите те, что стоят более 10\,000 рублей.  
    Кофемашина характеризуется брендом, моделью и ценой.  
    Дополнительно:
    \begin{itemize}
        \item если с молочным резервуаром — объём резервуара и возможность регулировки пены;
        \item если поддерживает капсулы — совместимые системы (Nespresso, Dolce Gusto и др.).
    \end{itemize}

    \item \textbf{Фотоаппараты.}  
    Введите информацию о фотоаппаратах и выведите те, у которых разрешение матрицы выше 20 МП.  
    Фотоаппарат характеризуется брендом, моделью и разрешением матрицы.  
    Дополнительно:
    \begin{itemize}
        \item если зеркальный — наличие оптического видоискателя и количество точек фокусировки;
        \item если поддерживает 4K-видео — максимальная частота кадров и битрейт.
    \end{itemize}

    \item \textbf{Часы.}  
    Введите данные о часах и выведите те, что стоят более 5000 рублей.  
    Часы характеризуются брендом, типом (механические/кварцевые/умные) и ценой.  
    Дополнительно:
    \begin{itemize}
        \item если водонепроницаемые — глубина погружения и стандарт (ISO 6425 и др.);
        \item если умные — поддержка ОС (Wear OS, watchOS) и время автономной работы.
    \end{itemize}

    \item \textbf{Книги.}  
    Введите информацию о книгах и выведите те, у которых количество страниц превышает 300.  
    Книга характеризуется названием, автором и количеством страниц.  
    Дополнительно:
    \begin{itemize}
        \item если учебник — предмет и класс;
        \item если художественная — жанр и год первого издания.
    \end{itemize}

    \item \textbf{Фильмы.}  
    Введите данные о фильмах и выведите те, чей рейтинг выше 7.0.  
    Фильм характеризуется названием, режиссёром и рейтингом.  
    Дополнительно:
    \begin{itemize}
        \item если научная фантастика — наличие сиквелов и бюджет;
        \item если анимационный — студия-производитель и возрастное ограничение.
    \end{itemize}

    \item \textbf{Рестораны.}  
    Введите информацию о ресторанах и выведите те, у которых средний чек ниже 1000 рублей.  
    Ресторан характеризуется названием, кухней и средним чеком.  
    Дополнительно:
    \begin{itemize}
        \item если доставка — минимальная сумма заказа и зона доставки;
        \item если есть детское меню — возрастные категории и наличие игровой зоны.
    \end{itemize}

    \item \textbf{Спортивные залы.}  
    Введите данные о спортивных залах и выведите те, у которых абонемент дешевле 3000 рублей/мес.  
    Зал характеризуется названием, адресом и стоимостью месячного абонемента.  
    Дополнительно:
    \begin{itemize}
        \item если есть бассейн — длина дорожки и температура воды;
        \item если круглосуточный — наличие охраны и системы видеонаблюдения.
    \end{itemize}

    \item \textbf{Музыкальные инструменты.}  
    Введите информацию об инструментах и выведите те, что стоят менее 10\,000 рублей.  
    Инструмент характеризуется типом (гитара, синтезатор и т.д.), брендом и ценой.  
    Дополнительно:
    \begin{itemize}
        \item если электронный — количество клавиш и наличие MIDI-выхода;
        \item если струнный — материал струн и наличие кейса.
    \end{itemize}
\end{enumerate}
